% Generated draft: requires mathematical and source review.
\documentclass[11pt]{article}
\usepackage[margin=25mm]{geometry}
\usepackage{fontspec}
\setmainfont{Latin Modern Roman}
\usepackage{amsmath,amssymb,mathtools,unicode-math}
\setmathfont{Latin Modern Math}
\usepackage{enumitem,needspace,xurl,hyperref}
\setlength{\parindent}{0pt}
\setlength{\parskip}{5pt}
\setlength{\emergencystretch}{3em}
\newcommand{\N}{\mathbb{N}}
\newcommand{\Z}{\mathbb{Z}}
\newcommand{\Q}{\mathbb{Q}}
\newcommand{\R}{\mathbb{R}}
\newcommand{\C}{\mathbb{C}}
\newcommand{\F}{\mathbb{F}}
\newcommand{\A}{\mathbb{A}}
\newcommand{\PP}{\mathbb{P}}
\newcommand{\E}{\mathbb{E}}
\newcommand{\eps}{\varepsilon}
\newcommand{\dd}{\,\mathrm{d}}
\begin{document}
\section*{UnsolvedMath ID: 10900062}
\textbf{Source problem code:} catalogue entry
\subsection{Definitions and mathematical statement}
Let \(L=L_1\cup L_2\subset S^3\) be a two-component link. A joining band is an embedded rectangle \(b:[0,1]^2\to S^3\) whose opposite attaching edges lie on small arcs of the two distinct components and whose interior is disjoint from \(L\). Replacing the attaching arcs by the other two edges produces a knot \(K_b\). Its core is the arc \(\beta=b([0,1]\times\{1/2\})\).

The source calls \(b\) complicated when either of the following holds:
\begin{enumerate}
\item When \(L\) is split, the core cannot be isotoped to meet a splitting sphere in fewer than three points. A splitting sphere is an embedded \(S^2\subset S^3\setminus L\) separating \(L_1\) from \(L_2\).
\item The core cannot be isotoped disjoint from any minimal-genus Seifert surface for \(L\).
\end{enumerate}
The alternatives are joined by ``or.'' The allowable core isotopies, fixed-endpoint conditions and quantification over splitting spheres require source confirmation; the first alternative is stated only for split \(L\). The orientation and connectedness conventions for minimal-genus Seifert surfaces are also absent.

Taylor's Problem 8.8 asks for an explicit hyperbolic knot \(K\), meaning one with complete finite-volume hyperbolic complement, such that no choice of two-component link and complicated joining band yields \(K\):
\[
 \exists K\text{ hyperbolic}\quad
 \forall(L,b):\quad b\text{ complicated}\ \Longrightarrow\ K_b\not\cong K.
\]
Thus ruling out one selected band presentation is insufficient; all such presentations must be excluded.

The source reports Taylor's result that knots produced using complicated bands satisfy the Cabling Conjecture. This does not say that such knots are themselves cable knots, despite the summary's compressed ``shown cabling'' language. The supplied 2026 assessment reports no located explicit example answering the stated exclusion problem; no independent search or resolution is asserted here.
\subsection{Short English statement}
Give a hyperbolic knot that has no presentation as a band sum of a two-component link using a complicated band. Complication is defined by the source's split-sphere or Seifert-surface obstruction, whose isotopy and quantifier conventions need checking. The exclusion must cover every presentation.
\subsection{Sources}
\begin{itemize}
\item Local UnsolvedMath registry, permanent ID \(10900062\), exact source code \texttt{catalogue entry}, subsection 8.8, attributed to Taylor; no proposed year supplied.
\item Delp, Hoffoss and Manning, \emph{Problems In Groups, Geometry, and Three-Manifolds} (2015), 8.8, PDF page 6: \url{https://arxiv.org/pdf/1512.04620#page=6}; supplied status-source URL \url{https://arxiv.org/abs/1512.04620}.
\item The source reports Taylor's Cabling-Conjecture result for complicated band attachments. The summary names Taylor, ``A note on band sums,'' and related work, without year or full bibliography; its claim of verification is a supplied assessment rather than a check performed here.
\item Source-TeX extraction checked against rendered PDF text, accessed 2026-07-29. Source evidence NEEDS\_REVIEW: open in 2015, no complete modern primary-literature audit recorded. Registry status open and classification OPEN-TRIAGE; summary reports no requested example located as of 2026.
\item Dataset attribution: Ulam AI / UnsolvedMath Contributors, CC BY 4.0. Source rights NEEDS\_REVIEW: publicly accessible author-submitted source TeX and PDF, redistribution terms requiring release review.
\end{itemize}
\end{document}
